\section{古典概型和几何概型}

上节介绍了概率的定义和运算法则，但没有给出概率的计算方法。
本节介绍古典概型和几何概型及其计算方法。
古典概型是概率论的雏形，也是最简单、最常见的类型。

本节要点：
\begin{itemize}
    \item 掌握古典概型和几何概型的概念；
    \item 熟练掌握两种概型的计算方法。
\end{itemize}

%============================================================
\subsection{排列组合的概念}

\begin{definition}[排列和组合]
从$n$个元素中任取$m$个元素（$m\leqslant n$）排成一有序列，称{\bf 从$n$个元素中取出$m$个元素的一个排列}，所有排列的个数称作{\bf 排列数}，记作$P_{n}^{m}$，有：
\[
P_{n}^{m}:=\frac{n!}{\left( n-m \right) !}
\]
从$n$个元素中任取$m$个元素（$m\leqslant n$）组成一无序组，称{\bf 从$n$个元素中取出$m$个元素的一个组合}，所有组合的个数称作{\bf 组合数}，记作$C_{n}^{m}$，有：
\[
C_{n}^{m}:=\frac{n!}{\left( n-m \right) !\cdot m!}=\frac{P_{n}^{m}}{m!}
\]
特别地，$P_{n}^{1}=C_{n}^{1}=n,P_{n}^{n}=n!,C_{n}^{n}=1$。
\end{definition}

排列组合均是对“从$n$个元素中抽取$m$个元素”的计算，排列突出有序，组合突出无序。
排列组合的关系还可以这么理解，当$m$个元素排列完成后，若不考虑其顺序，就是组合，由于每个排列有$P_{m}^{m}$种不同的顺序，所以：
\[
P_{n}^{m}=C_{n}^{m}\cdot P_{m}^{m}=C_{n}^{m}\cdot m!
\]

%============================================================
\subsection{古典概型的概念}

\begin{definition}[古典概型]
若随机试验$E$满足：
\begin{itemize}
    \item 其样本空间$\varOmega $中只有有限个样本点；
    \item 每个基本事件出现的概率相等，
\end{itemize}
则称试验$E$为{\bf 古典概型试验}，或{\bf 等可能概型试验}，上述条件称为{\bf 古典概型特征}。
对于任意事件$A$，其概率$P\left( A \right) $为：
\[
P\left( A \right) :=\frac{A\text{的样本数}}{\text{总样本数}}
\]
\end{definition}

古典概型的计算基于排列组合。
在古典概型中，每个样本点具有等概率$\frac{1}{n}$，而事件就是样本点的集合，有几个样本点就是$\frac{1}{n}$的几倍。

%============================================================
\subsection{几何概型的概念}

\begin{definition}[几何概型]
若随机试验$E$满足：
\begin{itemize}
    \item 其样本空间$\varOmega $为一几何区域（一维、二维或更高任意维）；
    \item 每个基本事件出现的概率相等，
\end{itemize}
则称试验$E$为{\bf 几何概型试验}，上述条件称为{\bf 几何概型特征}。
对于任意事件$A$，其概率$P\left( A \right) $为：
\[
P\left( A \right) :=\frac{A\text{的几何测度}}{\varOmega \text{的几何测度}}
\]
\end{definition}

古典概型中样本数量是有限的，几何概型中样本数量是无限的、不可数的或是稠密的，两者分别对应离散和连续的样本空间。
通俗来讲，古典概型的计算就是“掰手指”，而几何概型就是“量尺子”。

%============================================================
\subsection{例}

\begin{example}
设有4个不同的箱子，3个不同的球，每个球可以放入任意箱子，且概率相等，计算：
\begin{enumerate}
    \item 3个球放入不同的3个箱子的概率是多少，
    \item 1号箱子和2号箱子都有球的概率是多少。
\end{enumerate}
\end{example}

解：

1. 首先计算所有可能的放置方法，每个球有4种放法，3个球共有$4^3$种放法。
然后考虑3个球放入不同的3个箱子，第一个球有4个选择，第二球剩下3个箱子选择，最后一个球只剩下2个箱子选择，所以共有$4\cdot 3\cdot 2=24$种放法，于是概率：
\[
P=\frac{4\cdot 3\cdot 2}{4^3}=\frac{3}{8}
\]

2. 1号箱子和2号箱子都有球，可分为3种情况：
\begin{itemize}
    \item 1号箱子1个球，2号箱子1个球，剩下的球在3号或4号箱子，$C_{3}^{1}\cdot C_{2}^{1}\cdot 2=12$或$P_{3}^{2}\cdot 2=12$，$P\left( A_1 \right) =\frac{12}{4^3}=\frac{12}{64}$；
    \item 1号箱子1个球，2号箱子2个球，$C_{3}^{1}\cdot C_{2}^{2}=3$，$P\left( A_2 \right) =\frac{3}{4^3}=\frac{3}{64}$；
    \item 1号箱子2个球，2号箱子1个球，$C_{3}^{2}\cdot C_{1}^{1}=3$，$P\left( A_3 \right) =\frac{3}{4^3}=\frac{3}{64}$，
\end{itemize}
且这3种情况互斥，于是：
\[
P\left( A_1+A_2+A_3 \right) =P\left( A_1 \right) +P\left( A_2 \right) +P\left( A_3 \right) =\frac{9}{32}
\]

~

\begin{example}
将一根直棒随机分成3段，求3段能构成三角形的概率。
\end{example}

解：

假设棒长$a$，3段长度分别为$x,y,\left( a-x-y \right)$，于是样本空间为：
\[
\varOmega =\{\left( x,y \right) |0<x<a,0<y<a,a-x-y>0\}
\]
在{\it xy}平面上为{\it x}轴、{\it y}轴和直线$y=a-x$构成的面积。
能构成三角的充要条件是任意两边之和大于第三边，于是：
\[
\begin{cases}
	x+y>a-x-y\\
	x+\left( a-x-y \right) >y\\
	y+\left( a-x-y \right) >x\\
\end{cases}\qquad \Rightarrow \qquad \begin{cases}
	x+y>\frac{a}{2}\\
	y<\frac{a}{2}\\
	x<\frac{a}{2}\\
\end{cases}
\]
在\textit{xy}平面上为3条直线$x=\frac{a}{2}$、$y=\frac{a}{2}$和$y=\frac{a}{2}-x$构成的面积，可得概率为$1/4$。

\begin{figure}[h]
\centering
\includegraphics[height=3cm]{1.3.png}
\end{figure}




